\section{Preliminaries}
\label{sec-pre}

\stitle{RDF graphs and schemas}.
Let $I$, $B$, and $L$ be disjoint sets of IRIs, blank nodes, and
literals.
An RDF graph $\G$ is a finite set of triples
$(s, p, o) \in (I \cup B) \times I \times (I \cup B \cup L)$.
A schema $\Sc = (C, P, \dom, \rng, \sub)$ over $\G$ consists of a set
of classes $C \subseteq I$, a set of properties $P \subseteq I$, the
domains $\dom(p) \subseteq C$ and ranges
$\rng(p) \subseteq C \cup D$ of each property, where $D$ is a set of
datatypes, and the reflexive--transitive closure $\sub$ of
\kw{rdfs{:}subClassOf}~\cite{rdfs11}.
A type is a class or a datatype; $\Lit$ denotes the datatypes and
\kw{rdfs{:}Literal}.

\stitle{The SPARQL algebra}.
Let $\V$ be a set of variables, disjoint from $I \cup B \cup L$.
A triple pattern is a triple over $I \cup L \cup \V$.
Following Pérez et al.~\cite{perez2009semantics} and the SPARQL 1.1
specification~\cite{sparql11}, a solution mapping $\mu$ is a partial
function from $\V$ to $I \cup B \cup L$, two mappings are compatible
when they agree on their shared variables, and a graph pattern $P$
evaluates on $\G$ to a multiset of mappings $\eval{P}{\G}$.
Patterns are built from basic graph patterns (BGPs) by the operators
in Table~\ref{tab-ops}.
Evaluating SPARQL graph patterns is PSPACE-complete in
general~\cite{perez2009semantics}, so executing a plan to inspect it is
not free; this motivates checking what can be checked without
execution.

\begin{table}[t]
  \caption{SPARQL algebra operators and the \hermes actions that apply them.}
  \label{tab-ops}
  \small
  \begin{tabular}{@{}lll@{}}
    \toprule
    Operator & Meaning & Action \\
    \midrule
    BGP, Join & conjunction of triple patterns & \kw{triple} \\
    Filter & keep mappings satisfying $e$ & \kw{filter} \\
    LeftJoin & optional extension & \kw{optional} \\
    Union & either branch & \kw{union} \\
    Minus & remove compatible mappings & \kw{minus} \\
    Extend & bind $?v$ to $e$ & \kw{bind} \\
    Group, Aggregate & group, then aggregate & \kw{group}, \kw{aggregate} \\
    Filter on groups & HAVING & \kw{having} \\
    OrderBy & sort & \kw{order} \\
    Project, Distinct & restrict, deduplicate & \kw{project}, \kw{distinct} \\
    Slice & LIMIT, OFFSET & \kw{limit}, \kw{offset} \\
    ASK form & existence & \kw{ask} \\
    \bottomrule
  \end{tabular}
\end{table}

\stitle{Natural language to SPARQL}.
Given a question $q$, a graph $\G$, and its schema $\Sc$, the task is to
produce a query $Q$ whose answers $\eval{Q}{\G}$ match the intended
answers of $q$.
Following the Text2SPARQL evaluation~\cite{text2sparql}, answers are
compared as sets of values: every value bound to every projected
variable counts once, and quality is measured by set precision, recall,
and $F_1$ against the answers of a gold query.
